\chapter{Лабораторная работа №9. Конфигурирование кампусной сети.}
\label{ch:chapter 9}

\section{Цели}

Лабораторная работа помогает получить практические навыки по изучению
следующих тем:

\begin{itemize}

\item Общие концепции и архитектура кампусной сети

\item Наиболее распространенные сетевые технологии

\item Жизненный цикл кампусных сетей: планирование и проектирование, развертывание и внедрение, эксплуатация и техобслуживание, а также оптимизация

\item Процесс реализации проекта кампусной сети

\end{itemize}

\section{Топология сети}

В офисном здании необходимо создать сеть. Офисное здание имеет шесть этажей.
В настоящее время в эксплуатацию введены три этажа: холл с рецепцией на первом этаже,
административный отдел и кабинет генерального директора на втором этаже,
отдел исследований и разработок и отдел маркетинга на третьем этаже.
Аппаратный зал с основным оборудованием располагается на первом этаже.

\includegraphics{lb9_1}

\section{Сбор и анализ требований}

В ходе проектирования были собраны и проанализированы следующие требования:

\begin{itemize}
    \item \textbf{Бюджет проекта:} Бюджет ограничен. Проект необходимо реализовать с минимальными затратами.
    \item \textbf{Типы подключаемых терминалов:} Будет осуществляться развертывание как проводных, так и беспроводных терминалов.
    \item \textbf{Количество терминалов:} Первый этаж: 10 проводных и 100 беспроводных. Второй и третий этажи: 200 проводных и 50 беспроводных.
    \item \textbf{Режим управления сетью:} Унифицированное управление сетью с помощью SNMP.
    \item \textbf{Объем сетевого трафика:} Большая часть трафика — внутренний. Скорость проводного доступа — 100 Мбит/с.
    \item \textbf{Требования к доступности:} Необходимо реализовать резервирование компонентов сети уровня 3.
    \item \textbf{Требования к безопасности:} Необходимо реализовать контроль сетевого трафика.
    \item \textbf{Режим доступа в Интернет:} Статические IP-адреса.
\end{itemize}

\section{Планирование и проектирование}

\textbf{Задача 1. Выбор устройства и проектирование физической топологии.}

На основе анализа требований была разработана физическая топология сети,
включающая уровни доступа, агрегации, ядра и выходную область.
Коммутаторы уровня доступа (F1-ACC1, F2-ACC1, F2-ACC2, F2-ACC3, F3-ACC1, F3-ACC2, F3-ACC3)
подключены к коммутаторам агрегации (F2-AGG1, F3-AGG1), которые в свою очередь подключены
к ядерному коммутатору CORE1. CORE1 подключен к маршрутизатору для выхода в Интернет.

\includegraphics{lb9_1}

\textbf{Задача 2. Сетевое проектирование уровня 2.}

Было выполнено планирование VLAN для проводной и беспроводной сети:

\begin{verbatim}
[S1] vlan batch 100 105 205
[S1] vlan batch 101 102 106 206
[S1] vlan batch 103 104 107 207
[S1] vlan batch 201 202 203 204
\end{verbatim}

Результаты сведены в таблицу:

\begin{table}[h]
\small
\centering
\caption{Планирование VLAN}
\begin{tabular}{|l|l|}
\hline
\textbf{Идентификатор VLAN} & \textbf{Описание} \\ \hline
1 & VLAN управления устройствами уровня 2 на первом этаже \\ \hline
2 & VLAN управления устройствами уровня 2 на втором этаже \\ \hline
3 & VLAN управления устройствами уровня 2 на третьем этаже \\ \hline
100 & VLAN для серверов \\ \hline
101 & VLAN для офиса генерального директора \\ \hline
102 & VLAN для административного отдела \\ \hline
103 & VLAN для отдела маркетинга \\ \hline
104 & VLAN для отдела исследований и разработок (R\&D) \\ \hline
105 & VLAN для беспроводных терминалов на первом этаже \\ \hline
106 & VLAN для беспроводных терминалов на втором этаже \\ \hline
107 & VLAN для беспроводных терминалов на третьем этаже \\ \hline
201 & VLAN для связи между F2-AGG1 и CORE1 \\ \hline
202 & VLAN для связи между F3-AGG1 и CORE1 \\ \hline
203 & VLAN для связи между F2-AGG1 и F3-AGG1 \\ \hline
204 & VLAN для связи между CORE1 и маршрутизатором \\ \hline
205 & VLAN управления беспроводной сетью на первом этаже \\ \hline
206 & VLAN управления беспроводной сетью на втором этаже \\ \hline
207 & VLAN управления беспроводной сетью на третьем этаже \\ \hline
\end{tabular}
\end{table}

\textbf{Задача 3. Сетевое проектирование уровня 3.}

Спланированы IP-сети, методы назначения адресов и режимы маршрутизации.
Пример конфигурации для сети управления устройствами:

\begin{verbatim}
[R1] interface Vlanif1
[R1-Vlanif1] ip address 192.168.1.254 255.255.255.0
[R1-Vlanif1] quit
\end{verbatim}

В качестве протокола динамической маршрутизации выбран OSPF:

\begin{verbatim}
[CORE1] ospf 1
[CORE1-ospf-1] area 0
[CORE1-ospf-1-area-0.0.0.0] network 192.168.1.0 0.0.0.255
[CORE1-ospf-1-area-0.0.0.0] network 192.168.100.0 0.0.0.255
[CORE1-ospf-1-area-0.0.0.0] network 192.168.201.0 0.0.0.3
[CORE1-ospf-1-area-0.0.0.0] quit
[CORE1-ospf-1] quit
\end{verbatim}

\textbf{Задача 4. Проектирование WLAN.}

Спроектирована беспроводная сеть. Для каждого этажа создается свой SSID.
Используется политика безопасности WPA-WPA2+PSK+AES.

Настройка профиля безопасности:

\begin{verbatim}
[AC] wlan
[AC-wlan-view] security-profile name WLAN-F1
[AC-wlan-sec-prof-WLAN-F1] security wpa-wpa2 psk pass-phrase WLAN@Guest123 aes
[AC-wlan-sec-prof-WLAN-F1] quit
\end{verbatim}

Настройка профиля SSID:

\begin{verbatim}
[AC-wlan-view] ssid-profile name WLAN-F1
[AC-wlan-ssid-prof-WLAN-F1] ssid WLAN-F1
[AC-wlan-ssid-prof-WLAN-F1] quit
\end{verbatim}

\section{Реализация}

\textbf{Задача 1. Схема конфигурирования.}

На основании плана были составлены схемы конфигурирования для каждого устройства.

\textbf{Конфигурация на маршрутизаторе:}

\begin{verbatim}
<Router> system-view
[Router] sysname Router
[Router] interface GigabitEthernet0/0/0
[Router-GigabitEthernet0/0/0] ip address 1.1.1.1 24
[Router-GigabitEthernet0/0/0] quit
[Router] interface GigabitEthernet0/0/1
[Router-GigabitEthernet0/0/1] ip address 192.168.204.1 24
[Router-GigabitEthernet0/0/1] quit
[Router] acl 2000
[Router-acl-basic-2000] rule 5 permit source 192.168.105.0 0.0.0.255
[Router-acl-basic-2000] rule 10 permit source 192.168.106.0 0.0.0.255
[Router-acl-basic-2000] rule 15 permit source 192.168.107.0 0.0.0.255
[Router-acl-basic-2000] quit
[Router] nat address-group 1 1.1.1.2 1.1.1.10
[Router] interface GigabitEthernet0/0/0
[Router-GigabitEthernet0/0/0] nat outbound 2000 address-group 1
[Router-GigabitEthernet0/0/0] nat server protocol tcp global current-interface 8080 inside 192.168.100.1 www
[Router-GigabitEthernet0/0/0] quit
[Router] ospf 1
[Router-ospf-1] default-route-advertise always
[Router-ospf-1] area 0
[Router-ospf-1-area-0.0.0.0] network 192.168.204.0 0.0.0.3
[Router-ospf-1-area-0.0.0.0] quit
[Router-ospf-1] quit
[Router] ip route-static 0.0.0.0 0 1.1.1.254
\end{verbatim}

\textbf{Конфигурация на CORE1:}

\begin{verbatim}
<CORE1> system-view
[CORE1] sysname CORE1
[CORE1] vlan batch 100 105 201 202 204 205
[CORE1] dhcp enable
[CORE1] ip pool sta-f1
[CORE1-ip-pool-sta-f1] gateway-list 192.168.105.254
[CORE1-ip-pool-sta-f1] network 192.168.105.0 mask 24
[CORE1-ip-pool-sta-f1] quit
[CORE1] interface Vlanif105
[CORE1-Vlanif105] ip address 192.168.105.254 24
[CORE1-Vlanif105] dhcp select global
[CORE1-Vlanif105] quit
[CORE1] interface Vlanif201
[CORE1-Vlanif201] ip address 192.168.201.1 30
[CORE1-Vlanif201] quit
[CORE1] interface GigabitEthernet0/0/2
[CORE1-GigabitEthernet0/0/2] port link-type access
[CORE1-GigabitEthernet0/0/2] port default vlan 201
[CORE1-GigabitEthernet0/0/2] quit
[CORE1] ospf 1
[CORE1-ospf-1] area 0
[CORE1-ospf-1-area-0.0.0.0] network 192.168.1.0 0.0.0.255
[CORE1-ospf-1-area-0.0.0.0] network 192.168.100.0 0.0.0.255
[CORE1-ospf-1-area-0.0.0.0] network 192.168.105.0 0.0.0.255
[CORE1-ospf-1-area-0.0.0.0] network 192.168.201.0 0.0.0.3
[CORE1-ospf-1-area-0.0.0.0] quit
[CORE1-ospf-1] quit
\end{verbatim}

\textbf{Задача 2. Приемка проекта.}

После завершения настройки были определены следующие пункты для приемки:

\begin{enumerate}
    \item Проверка возможности беспроводных клиентов обнаруживать беспроводные сигналы и успешно получать доступ к сети.
    \item Проверка успешного установления отношений соседства OSPF.
    \item Проверка возможности подключения внутри сетей и между сетями.
    \item Проверка контроля доступа для гостей, использующих беспроводные терминалы.
    \item Проверка возможности NMS управлять сетевыми устройствами.
\end{enumerate}

Проверка связности выполняется командой ping:

\begin{verbatim}
<R1> ping 192.168.103.1
  PING 192.168.103.1: 56  data bytes, press CTRL_C to break
  Reply from 192.168.103.1: bytes=56 Sequence=1 ttl=254 time=40 ms
  Reply from 192.168.103.1: bytes=56 Sequence=2 ttl=254 time=30 ms
  Reply from 192.168.103.1: bytes=56 Sequence=3 ttl=254 time=30 ms
  --- 192.168.103.1 ping statistics ---
  3 packet(s) transmitted
  3 packet(s) received
  0.00% packet loss
\end{verbatim}

\section{Эксплуатация и техническое обслуживание сети (O\&M)}

\textbf{Задача 1. Организация O\&M.}

После сдачи проекта были организованы мероприятия по техническому обслуживанию:

\begin{table}[h]
\tiny
\centering
\caption{План технического обслуживания}
\begin{tabular}{|l|l|l|}
\hline
\textbf{Периодичность} & \textbf{Пункт проверки} & \textbf{Критерий оценки} \\ \hline
Каждый день & Загрузка ЦП & Загрузка ЦП каждого модуля в норме (менее 80\%) \\ \hline
Каждый день & Аварийная информация & Аварийные сигналы анализируются и обрабатываются \\ \hline
Один раз в неделю & Температура окружающей среды & Диапазон от 0°C до 50°C \\ \hline
Один раз в месяц & Резервное копирование конфигурации & Копии конфигураций создаются ежемесячно \\ \hline
Один раз в месяц & Смена пароля & Смена паролей для входа в устройство \\ \hline
\end{tabular}
\end{table}

Пример проверки состояния устройства:

\begin{verbatim}
<CORE1> display cpu-usage
CPU Usage Stat. Cycle: 60 (Second)
CPU Usage      : 15% Max: 25%
CPU Usage Stat. Time : 2024-01-15  15:30:00

<CORE1> display alarm urgent
--------------------
No alarm information.
--------------------
\end{verbatim}

\section{Оптимизация сети}

\textbf{Задача 1. Оптимизация производительности.}

По мере развития предприятия внутренний трафик между вторым и третьим этажами существенно увеличился. Для оптимизации канала предложены следующие решения:

\begin{enumerate}
    \item Добавить физические каналы между F2-AGG1 и F3-AGG1 и настроить агрегирование каналов Ethernet (Eth-Trunk).
    \item Изменить значения стоимости OSPF и реализовать балансировку нагрузки, чтобы часть трафика передавалась через CORE1.
\end{enumerate}

Пример настройки агрегирования каналов:

\begin{verbatim}
[F2-AGG1] interface Eth-Trunk 1
[F2-AGG1-Eth-Trunk1] mode lacp
[F2-AGG1-Eth-Trunk1] quit
[F2-AGG1] interface GigabitEthernet0/0/3
[F2-AGG1-GigabitEthernet0/0/3] eth-trunk 1
[F2-AGG1-GigabitEthernet0/0/3] quit
[F2-AGG1] interface GigabitEthernet0/0/4
[F2-AGG1-GigabitEthernet0/0/4] eth-trunk 1
[F2-AGG1-GigabitEthernet0/0/4] quit
\end{verbatim}

\section{Вопросы}

\textbf{1. В этом проекте CORE1, F2-AGG1 и F3-AGG1 образуют физическое кольцо. Однако на этапе планирования и проектирования сети каналы, по которым осуществляется взаимодействие этих трех устройств, были назначены в разные VLAN. Следовательно, петля не формируется. Однако во время лабораторной работы вы обнаружили, что между двумя устройствами отношения соседства не устанавливаются или устанавливаются некорректно. Найдите основную причину и решение.}

\textbf{Ответ:}

\textbf{Основная причина:} Несмотря на то, что на уровне VLAN петли отсутствуют (трафик разных VLAN логически разделен), физически кольцевая топология сохраняется. Протокол STP (Spanning Tree Protocol) работает на физическом уровне и не учитывает принадлежность кадров к разным VLAN при передаче BPDU (Bridge Protocol Data Units). BPDU передаются без тегов VLAN. Поэтому, даже если каналы назначены в разные VLAN, STP все равно обнаруживает физическую петлю и блокирует один из портов (обытельно порт с наивысшей стоимостью или наименьшим приоритетом) для предотвращения широковещательного шторма. В результате один из линков между коммутаторами оказывается заблокированным протоколом STP, из-за чего отношения соседства OSPF между двумя устройствами не могут установиться (OSPF-пакеты не проходят через заблокированный порт).

\textbf{Решение:} Необходимо отключить STP на интерфейсах между CORE1, F2-AGG1 и F3-AGG1, так как петли на уровне 2 исключены на этапе проектирования (разные VLAN). Для этого выполняется команда:

\begin{verbatim}
[CORE1] interface GigabitEthernet0/0/2
[CORE1-GigabitEthernet0/0/2] stp disable
[CORE1-GigabitEthernet0/0/2] quit
\end{verbatim}

Аналогичные команды необходимо выполнить на соответствующих интерфейсах F2-AGG1 и F3-AGG1. После этого BPDU не будут передаваться, и STP не будет блокировать порты, что позволит OSPF корректно установить соседские отношения на всех трех линках.

\textbf{2. Что нового вы узнали благодаря этой лабораторной работе? Как эти знания смогут помочь вам в дальнейшей учебе или работе?}

\textbf{Ответ:}

Благодаря этой лабораторной работе я узнал(а) следующее:

\begin{enumerate}
    \item \textbf{Комплексный подход к проектированию сети:} Понял(а), что построение корпоративной сети начинается не с настройки оборудования, а со сбора и анализа требований заказчика (бюджет, количество пользователей, типы терминалов, требования к безопасности и доступности).
    \item \textbf{Иерархическая модель сети:} Освоил(а) принципы построения трехуровневой модели сети (ядро, агрегация, доступ) и明白了 назначение каждого уровня.
    \item \textbf{Планирование VLAN и IP-адресации:} Научился(ась) логически разделять сеть на сегменты (VLAN) для разных отделов и типов трафика (проводной/беспроводной), а также планировать IP-сети и методы назначения адресов (статический/DHCP).
    \item \textbf{Интеграция технологий:} Увидел(а), как различные технологии (VLAN, STP, OSPF, DHCP, NAT, WLAN, SNMP) объединяются в единое рабочее решение.
    \item \textbf{Жизненный цикл сети:} Изучил(а) полный цикл проекта: от планирования и проектирования до реализации, эксплуатации (мониторинг, резервное копирование) и последующей оптимизации.
\end{enumerate}

\textbf{Применение знаний:}

Эти знания помогут мне в дальнейшей учебе при изучении более сложных тем, таких как VPN, MPLS, SD-WAN, а также при выполнении курсовых и дипломных проектов, связанных с разработкой сетевой инфраструктуры. В профессиональной деятельности эти навыки позволят мне участвовать в реальных проектах по развертыванию сетей малого и среднего бизнеса, понимать логику существующих сетевых решений и предлагать обоснованные варианты их модернизации.